\section{Main Theorem}\label{sec:main}

\begin{theorem}[Conditional positive limiting loss]\label{thm:main}
Fix $\kappa\geq1$ and $q>0$.  Let
$r_0,C_{\mathrm{dim}},\delta,L_P,\zeta,C_T$ be positive quantities that
may depend only on $(\kappa,q)$, with $r_0\in\mathbb N$ and
\[
L_P<\delta/4,
\qquad
\zeta<\delta/4.
\]
Define
\[
\epsilon=
\left(\frac{\delta-L_P-\zeta}{\kappa^6C_T}\right)^2>0.
\]
For every $r\geq r_0$, every $n,k$ satisfying
Assumptions~\ref{assump:dimension} and \ref{assump:rank_window}, every
deterministic base triple satisfying
Assumption~\ref{assump:base_conditioning}, and the half-relaxed parallel ALS
trajectory under Assumptions~\ref{assump:gaussian_smoothing} and
\ref{assump:independent_initialization}, one has the event inclusion
\[
\mathsf C_2(\delta,L_P,\zeta,C_T)
\subseteq
\left\{
\begin{array}{l}
\displaystyle \lim_{t\to\infty}\mathcal L(X_t,Y_t,Z_t)
\text{ exists and is finite},\\[2mm]
\displaystyle \lim_{t\to\infty}\mathcal L(X_t,Y_t,Z_t)
\geq\epsilon\|T\|_F^2
\end{array}
\right\}.
\]
This is a deterministic implication on each realization in the displayed
event.  It makes no assertion that $\mathsf C_2$ is nonempty and no assertion
that it has positive probability.

There are no hidden constants in the displayed loss fraction.  Its exposed
dependence is exactly on $\kappa,\delta,L_P,\zeta,C_T$; $r,n,k$ enter only
through the theorem scope and certificate objects.  The probability mode is
event inclusion under the joint law, the residual floor used in the proof is
all-time, the conclusion is asymptotic, and all residuals and objectives use
the ambient Frobenius norm.

In the exact/noiseless coordinate-orthonormal specialization
$Q=I$, $E_\rho=0$, and $T=D_r$, the same entry, projector-path, and
finite-variation clauses give the stronger baseline conclusion
\[
\lim_{t\to\infty}\mathcal L(X_t,Y_t,Z_t)
\geq(\delta-L_P)^2\|T\|_F^2.
\]
In particular, because the conditional implication holds for every positive
choice with the stated margins, it implies the existential constant statement
in the formalized goal, with all theorem constants depending only on
$(\kappa,q)$.
\end{theorem}

The unresolved source-level problem is to prove a uniform bound
$\mathbb P[\mathsf C_2(\delta,L_P,\zeta,C_T)]\geq p_0(\kappa,q)>0$ for
some admissible constants.  Theorem~\ref{thm:main} neither proves nor assumes
such a bound.
